\section{Two-term dg kernels}\label{sec:realisation}

\Cref{thm:ambient} will show that profiles can be cut by detectors: once a kernel dg algebra \(B\) has one realisation with full profile and projective localisation, every far-away orthogonal \(\cU_E\) is realised.  The substantive question is which dg algebras admit such a realisation at all.  This section gives an explicit realisation for every proper dg algebra with cohomology in degrees \(-1,0\), and proves that these are exactly the proper connective dg algebras admitting a standard t-exact marked kernel realisation.  The converse holds for ordinary homological ring epimorphisms without target projectivity.  Full profile alone has a wider scope, as \Cref{ex:marking-boundary} shows.

\begin{question}[Kernel realisation]\label{q:realisation}
Which proper cohomologically non-positive dg algebras \(B\) are, up to a specified dg Morita identification, kernel dg algebras of stably flat universal localisations of finite-dimensional algebras at finite \(\Sigma\) with full bounded profile?
\end{question}

By \Cref{thm:heart-amplitude}, for a silting marking full profile is equivalent to fixed finite Tor-amplitude of the heart bimodule; it is a property of the embedding, not of \(B\) alone.

\subsection{Necessary conditions}

\begin{proposition}[Thin generators]\label{prop:thin}
In every kernel as in \Cref{q:realisation}, the transported sum \(C_0\) of the presentation cones is a classical generator of \(\per B\) with \(\Hom(C_0,C_0[n])=0\) for \(n\notin[-1,1]\).  If \(C_0\) is silting, \(H^n\RHom(C_0,C_0)=0\) for \(n\notin[-1,0]\).
\end{proposition}

\begin{proof}
The ambient representative of \(C_0\) is a two-term complex of finite projectives in degrees \(-1,0\); its mapping complex has terms only in degrees \(-1,0,1\) and computes derived Hom.  The cones classically generate the compact kernel by \Cref{lem:compact-intersection}.
\end{proof}

The statement concerns the existence of a generator with this property; a silting generator obtained by shifting summands, as in \Cref{sec:nakayama}, need not have two-term amplitude.  The silting hypothesis cannot be dropped: for \(R=k\times k\), \(p=(1,0)\), the universal localisation at the zero map \(Rp\to Rp\) is the projection onto the second factor, stably flat with projective target, and its cone \(Rp[1]\oplus Rp\) has \(H^1\) of its endomorphism dg algebra of dimension one.

\begin{proposition}[Higher multiplications]\label{prop:higher-mult}
Let \(H\) be a minimal \(A_\infty\)-algebra with underlying graded space \(H^{-1}\oplus H^0\).  Then the only possible higher multiplication is
\[
 m_3:(H^0)^{\otimes3}\longrightarrow H^{-1},
\]
and in a strictly unital model it is a normalised Hochschild \(3\)-cocycle of the algebra \(\Gam=H^0\) with values in the bimodule \(U=H^{-1}\).
\end{proposition}

\begin{proof}
The operation \(m_n\) has degree \(2-n\).  For \(n\ge4\) all inputs have non-positive degree and the output has degree at most \(-2\), where \(H\) vanishes; for \(n=3\) a negative input forces the output below \(-1\).  The \(A_\infty\) relation in arity \(3\) makes \(m_2\) associative, and the relation in arity \(4\) on degree-zero inputs is exactly the Hochschild cocycle condition; compositions of two \(m_3\) vanish for degree reasons.  Minimal models exist over a field by the transfer theorem, see \cite[Section~3.3]{Keller01Ainf}.
\end{proof}

With the graded algebra and identity linear marking fixed, the class of \(m_3\) is the remaining invariant of such a structure.

\begin{lemma}[Gauge classification]\label{lem:gauge}
Two strictly unital minimal \(A_\infty\)-structures on \(\Gam\oplus U[1]\) with the same associative product on \(\Gam\) and the same bimodule structure on \(U\) are \(A_\infty\)-isomorphic by an isomorphism with identity linear part if and only if their \(m_3\) differ by a Hochschild coboundary.  Thus such structures are classified by \(HH^3(\Gam,U)\).
\end{lemma}

\begin{proof}
An \(A_\infty\)-morphism \(f\) with \(f_1=\id\) has components \(f_n\) of degree \(1-n\); on \(\Gam\oplus U[1]\) the only possibly nonzero higher component is \(f_2:\Gam^{\otimes2}\to U\).  The arity-\(2\) morphism equation is vacuous since \(m_1=0\), so \(m_2\) is unchanged; the arity-\(3\) equation reads \(m_3'-m_3=\pm d_{\mathrm{Hoch}}(f_2)\); every higher equation involves a composite with two negative-degree outputs and holds automatically.  Conversely any normalised \(2\)-cochain \(f_2\) defines such an isomorphism.
\end{proof}

\subsection{Two-term dg algebras as crossed bimodules}

\begin{lemma}[Two-term dg algebras]\label{lem:crossed}
A unital dg algebra \(D\) with \(D^i=0\) for \(i\notin\{-1,0\}\) is the same as an algebra \(A=D^0\), an \(A\)-bimodule \(M=D^{-1}\), and a bimodule map \(\delta:M\to A\) satisfying
\begin{equation}\label{eq:peiffer}
 \delta(m)m'=m\delta(m')\qquad(m,m'\in M);
\end{equation}
the product of two elements of \(M\) is zero.  Put \(\Gam=A/\delta(M)\) and \(U=\ker\delta\).  Then \(\delta(M)\) annihilates \(U\) on both sides, \(U\) is a \(\Gam\)-bimodule, \(H^0(D)=\Gam\), \(H^{-1}(D)=U\), and
\begin{equation}\label{eq:crossed-extension}
 0\longrightarrow U\longrightarrow M\xrightarrow{\ \delta\ }A\longrightarrow\Gam\longrightarrow0
\end{equation}
is a crossed extension of \(\Gam\) by \(U\) in the sense of Baues--Minian \cite{BauesMinian02}.
\end{lemma}

\begin{proof}
The products in degrees \((0,0)\), \((0,-1)\), \((-1,0)\) give the algebra and bimodule structures, and products in degree \(-2\) vanish.  The Leibniz rule with one factor of degree zero says that \(\delta\) is a bimodule map; applied to the zero product \(mm'\) it gives \(0=\delta(m)m'-m\delta(m')\).  Conversely these identities are the Leibniz rule in every degree pair, and the bimodule axioms give associativity.  For \(u\in U\) and \(m\in M\), \eqref{eq:peiffer} gives \(\delta(m)u=m\delta(u)=0\) and \(u\delta(m)=\delta(u)m=0\).
\end{proof}

\begin{lemma}[The minimal model of a two-term dg algebra]\label{lem:two-term-minimal}
Let \(D=[M\xrightarrow{\delta}A]\) be as in \Cref{lem:crossed}, with \(\Gam\), \(U\) finite-dimensional.  Choose a \(k\)-linear section \(s:\Gam\to A\) of the projection with \(s(1)=1\), a complement \(M=U\oplus M'\), and for \(x,y\in\Gam\) let \(\beta(x,y)\in M'\) be the unique element with
\begin{equation}\label{eq:beta}
 s(x)s(y)-s(xy)=\delta\beta(x,y).
\end{equation}
Then the transferred minimal \(A_\infty\)-structure on \(\Gam\oplus U[1]\) has \(m_2\) the product of \(\Gam\) and the bimodule structure of \(U\), and
\begin{equation}\label{eq:m3-cocycle}
 m_3(x,y,z)=x\beta(y,z)-\beta(xy,z)+\beta(x,yz)-\beta(x,y)z\in U,
\end{equation}
the characteristic cocycle of the crossed extension \eqref{eq:crossed-extension}.  Its class in \(HH^3(\Gam,U)\) is independent of all choices and is the characteristic class of \eqref{eq:crossed-extension}.
\end{lemma}

\begin{proof}
Let \(i:\Gam\oplus U[1]\to D\) be \(s\) on \(\Gam\) and the inclusion on \(U\); let \(p:D\to\Gam\oplus U[1]\) be the projection \(A\to\Gam\) on \(A\) and the projection to \(U\) along \(M'\) on \(M\); let \(h:D\to D[-1]\) vanish on \(s(\Gam)\) and on \(M\), and on \(\delta(M)\) let \(h(\delta m')=-m'\) for \(m'\in M'\).  Then \(pi=\id\) and \(ip-\id=dh+hd\): on \(a=s(x)+\delta(m')\) both sides equal \(-\delta(m')\), and on \(m=u+m'\) both sides equal \(-m'\).  The transfer formula gives
\[
 m_3(x,y,z)=p\bigl(h(s(x)s(y))\,s(z)\bigr)-p\bigl(s(x)\,h(s(y)s(z))\bigr)
 =p\bigl(s(x)\beta(y,z)-\beta(x,y)s(z)\bigr),
\]
using \(h(s(x)s(y))=h(s(xy))+h(\delta\beta(x,y))=-\beta(x,y)\).  Applying \(\delta\) to the right-hand side of \eqref{eq:m3-cocycle} and expanding with \eqref{eq:beta} gives zero, so that cocycle lies in \(U\); hence the \(M'\)-component of \(s(x)\beta(y,z)-\beta(x,y)s(z)\) equals \(\beta(xy,z)-\beta(x,yz)\), and projecting to \(U\) along \(M'\) yields \eqref{eq:m3-cocycle}.  Different choices of \(s\), \(M'\), or of the transfer data change the minimal structure by an \(A_\infty\)-isomorphism with identity linear part \cite[Section~3.3]{Keller01Ainf}, hence \(m_3\) by a coboundary (\Cref{lem:gauge}).  That \eqref{eq:m3-cocycle} is the characteristic cocycle of the crossed extension is the construction of \cite[Section~3]{BauesMinian02}.
\end{proof}

\subsection{Realisation of a strict two-term dg algebra}

\begin{theorem}[Realisation of two-term dg algebras]\label{thm:realisation}
Let \(D=[M\xrightarrow{\delta}A]\) be a finite-dimensional unital dg algebra concentrated in degrees \(-1,0\), and let
\begin{equation}\label{eq:R-delta}
 R_\delta=\begin{pmatrix}A&A\\M&A\end{pmatrix},
 \qquad
 (a,v,m,b)(a',v',m',b')=\bigl(aa'+v\delta(m'),\ av'+vb',\ ma'+bm',\ bb'+\delta(m)v'\bigr),
\end{equation}
the Morita context ring of the pairings \(v\otimes m\mapsto v\delta(m)\) and \(m\otimes v\mapsto\delta(m)v\).  Let \(p=e_{11}\), \(e=e_{22}\), let \(s\) be the upper-right copy of \(1_A\), and \(\sigma:R_\delta p\to R_\delta e\), \(x\mapsto xs\).  Then:
\begin{enumerate}[label=\textup{(\roman*)}]
\item\label{re:i} the ring map \(\lambda_\delta:R_\delta\to S=M_2(A)\), \((a,v,m,b)\mapsto\begin{pmatrix}a&v\\\delta(m)&b\end{pmatrix}\), is the universal localisation at \(\sigma\);
\item\label{re:ii} \(S_{R_\delta}\cong(pR_\delta)^{\oplus2}\) and \({}_{R_\delta}S\cong(R_\delta e)^{\oplus2}\); in particular \(\lambda_\delta\) is stably flat;
\item\label{re:iii} for \(T=\Cone(\sigma)\) there is a strict quasi-isomorphism of dg algebras \(D\to\End^\bullet_{R_\delta}(T)^{\op}\);
\item\label{re:iv} the Morita functor \(F=T\Ltensor_D-:D(D)\xrightarrow{\sim}\cK_{\lambda_\delta}\) is t-exact for the standard t-structures, with \(H^n(FY)\cong(0,H^nY)\) as \(R_\delta\)-modules through \(R_\delta\to\Gam\); it preserves and reflects finite-dimensional total cohomology, so the bounded profile is all of \(\Dfd(D)\); and the heart bimodule is \(T\Ltensor_D\Gam\simeq h\Gam\), where \(hN=(0,N)\) denotes the \(R_\delta\)-module supported on the lower-right corner, of Tor-amplitude \([0,0]\).
\end{enumerate}
Moreover \(\dim_kR_\delta=3\dim_kA+\dim_kM\) and \(\dim_kS=4\dim_kA\).
\end{theorem}

\begin{proof}
\emph{Associativity and the ring map.}  The diagonal products are products in \(A\); each off-diagonal associativity identity is one of the bimodule identities for the regular bimodule \(A\) or for \(M\), except the identity \(\bigl((m\cdot v')\cdot m''\bigr)=m\cdot(v'\cdot m'')\), which reads \(\delta(m)v'm''=mv'\delta(m'')\) and follows from \eqref{eq:peiffer} applied to \(m\) and \(v'm''\) since \(\delta\) is a bimodule map; and the identity \((v\cdot m')\cdot v''=v\cdot(m'\cdot v'')\), which is associativity in \(A\).  The unit is \((1,0,0,1)\).  The map \(\lambda_\delta\) is multiplicative because \(\delta(ma'+bm')=\delta(m)a'+b\delta(m')\); its kernel is \(\{(0,0,u,0):u\in U\}\).

\ref{re:i}.  Since \(\lambda_\delta(s)=E_{12}\) has the two-sided inverse \(E_{21}\) between the corner projectives of \(M_2(A)\), \(\lambda_\delta\) inverts \(\sigma\).  Let \(f:R_\delta\to D'\) invert \(\sigma\), put \(p'=f(p)\), \(e'=f(e)\), \(s'=f(s)\), and let \(t\in e'D'p'\) satisfy \(s't=p'\), \(ts'=e'\).  For \(m\in M\) one has \(ms=(0,0,0,\delta(m))\), so \(f(m)=f(m)p'=f(m)s't=f(\delta(m)_e)\,t\), where \(\gamma_e\) denotes \(\gamma\in A\) in the lower-right corner; thus \(f\) is determined on \(M\) by its restriction to the diagonal and by \(t\).  Also \(\gamma_ps=s\gamma_e\) for \(\gamma\in A\), so \(f(\gamma_p)s'=s'f(\gamma_e)\) and the two diagonal copies of \(A\) are identified through \(s'\) and \(t\).  The elements \(p',e',s',t\) are matrix units, and \(\gamma\mapsto f(\gamma_p)+f(\gamma_e)\), \(E_{21}\mapsto t\) defines a ring map \(M_2(A)\to D'\).  It extends \(f\): on \(M\) it sends \(\lambda_\delta(m)=\delta(m)E_{21}\) to \(f(\delta(m)_e)t=f(m)\).  Uniqueness holds because \(t\) is the unique inverse of the corner isomorphism.

\ref{re:ii}.  The right module \(pR_\delta\) is the row \((A,A)\) with action \((x,y)(a,v,m,b)=(xa+y\delta(m),\ xv+yb)\), and each row of \(M_2(A)\) restricted along \(\lambda_\delta\) has this action, so \(S_{R_\delta}\cong(pR_\delta)^{\oplus2}\).  The left module \(R_\delta e\) is the column \((A,A)^{\mathsf T}\) with action \((a,v,m,b)\binom xy=\binom{ax+vy}{\delta(m)x+by}\), and each column of \(S\) has this action, so \({}_{R_\delta}S\cong(R_\delta e)^{\oplus2}\).  A ring epimorphism with projective target on one side has \(\Tor^{R_\delta}_i(S,S)=0\) for \(i>0\); universal localisations are ring epimorphisms, so \(\lambda_\delta\) is stably flat.

\ref{re:iii}.  The endomorphism complex of \(T=[R_\delta p\xrightarrow{s}R_\delta e]\), written with coefficients of right multiplications, is
\[
 M\xrightarrow{\ m\mapsto(\delta(m),\delta(m))\ }A\oplus A\xrightarrow{\ (a,b)\mapsto a-b\ }A
\]
in degrees \(-1,0,1\): the degree \(-1\) term is \(\Hom_{R_\delta}(R_\delta e,R_\delta p)=eR_\delta p=M\), and its differential consists of the two composites \(x\mapsto xsm\) and \(x\mapsto xms\), which are right multiplications by \(sm=\delta(m)_p\) and \(ms=\delta(m)_e\); the degree \(0\) term is \(pR_\delta p\oplus eR_\delta e\), with differential the commutator with \(s\); the degree \(1\) term is \(pR_\delta e=A\).  Its cohomology is \(U\) in degree \(-1\), \(\{(a,a)\}/\{(\delta m,\delta m)\}\cong\Gam\) in degree \(0\), and \(0\) in degree \(1\).  Send \(a\in A\) to simultaneous right multiplication by \(a\) on \(R_\delta p\) and \(R_\delta e\), and \(m\in M\) to right multiplication by \(m\) as a map \(R_\delta e\to R_\delta p\).  This is compatible with the differentials, \(d(a,a)=0\) and \(d(m)=(\delta m,\delta m)\), and multiplicative in the opposite algebra: composing ``right \(a\), then right \(m\)'' is right multiplication by \(am\), so \(a\cdot^{\op}m=am\) and \(m\cdot^{\op}a=ma\) with the corner at \(e\) acting on the left and the corner at \(p\) on the right, which is why the opposite is required; products of two degree \(-1\) maps vanish because \(T\) has no degree \(-2\) endomorphisms.  The map induces the isomorphisms \(\Gam\to H^0\) and \(U\to H^{-1}\) just computed, so it is a quasi-isomorphism.  By stable flatness, \(\cK_{\lambda_\delta}=\Loc(T)\) and \(\Kc_{\lambda_\delta}=\thick(T)\) (\Cref{lem:compact-intersection}), and dg Morita theory gives \(F\).

\ref{re:iv}.  Restrict \(T\), with its right \(D\)-action through \ref{re:iii}, to the two idempotents.  The \(e\)-component is \(eT=[eR_\delta p\to eR_\delta e]=[M\xrightarrow{\delta}A]=D\) as a right dg \(D\)-module.  The \(p\)-component is \(pT=[A\xrightarrow{1}A]\), on which \(A\) acts by right multiplication and \(m\in M\) acts from degree \(0\) to degree \(-1\) by multiplication by \(\delta(m)\); the degree \(-1\) map \(h\) equal to the identity on \((pT)^0=A\) and zero on \((pT)^{-1}\) is right \(D\)-linear, since on degree \(-1\) elements both \(h(x\cdot m)\) and \(h(x)\cdot m\) vanish for degree reasons, and satisfies \(dh+hd=1\).  Thus \(pT\) is contractible as a right dg \(D\)-module.  For a semi-free \(Y\), \(FY=T\otimes_DY\) has \(e\)-component \(D\otimes_DY=Y\) and contractible \(p\)-component, so \(H^n(FY)\cong(0,H^nY)\) with \(R_\delta\) acting through \(R_\delta\to R_\delta/R_\delta pR_\delta=\Gam\).  This proves t-exactness and the profile assertion.  For \(Y=\Gam\) all identifications commute with the regular right \(\Gam\)-action, giving \(T\Ltensor_D\Gam\simeq h\Gam\).  Independently, \ref{re:ii} and \Cref{thm:kernel-dictionary} give full profile.  The dimension count is immediate.
\end{proof}

\begin{remark}\label{rem:realisation-benchmarks}
For \(\delta=0\), \eqref{eq:R-delta} is the Morita context ring with zero pairings of \cite{GreenPsaroudakis14,GaoPsaroudakis17}, and the theorem realises the formal dg algebra \(A\oplus M[1]\); for \(M=0\) it is the localisation \(T_2(A)\to M_2(A)\) of the upper triangular ring, with ordinary kernel \(A\); if \(\delta\) is an isomorphism the kernel is zero.  Matrix-ring universal localisations of exact contexts and their homologicality criterion are due to Chen--Xi \cite{ChenXi19}; the content here is the identification of the kernel of the cone with the prescribed two-term dg algebra, differential included.
\end{remark}

\subsection{Every thin dg algebra occurs}

Finite strict models for proper connective dg algebras are known over arbitrary fields by Raedschelders--Stevenson \cite[Appendix~A, Theorem~A.3]{Goodbody24}; see \Cref{rem:finite-strict-priority} below.  The next construction gives such a two-term model directly from a prescribed Hochschild class.

\begin{theorem}[Finite crossed models]\label{thm:finite-crossed}
Let \(\Gam\) be a finite-dimensional algebra, \(U\) a finite-dimensional \(\Gam\)-bimodule, and \(\mu\in HH^3(\Gam,U)\).  There is a finite-dimensional two-term dg algebra \(D=[M\to A]\) with \(H^0(D)=\Gam\), \(H^{-1}(D)=U\) and minimal-model class \([m_3]=\mu\).  Consequently every proper dg algebra \(B\) with \(H^i(B)=0\) for \(i\notin\{-1,0\}\) is quasi-isomorphic to a finite-dimensional two-term dg algebra, and hence, by \Cref{thm:realisation}, is the kernel dg algebra of a stably flat universal localisation \(R\to M_2(A)\) of finite-dimensional algebras at a single map, with target projective on both sides, t-exact kernel embedding, and full bounded profile.
\end{theorem}

\begin{proof}
Let \(\Gam^e=\Gam\otimes_k\Gam^{\op}\) and let \(I=\Hom_k(\Gam^e,U)\) with \((qf)(r)=f(rq)\); it is a finite-dimensional injective left \(\Gam^e\)-module, being coinduced from \(k\), and \(\nu:U\to I\), \(\nu(u)(r)=ru\), is an injective \(\Gam^e\)-linear map since \(\nu(u)(1)=u\).  Put \(J=I/U\) with quotient map \(q:I\to J\).  Since \(HH^n(\Gam,I)=\Ext^n_{\Gam^e}(\Gam,I)=0\) for \(n\ge1\), the long exact sequence gives an isomorphism \(\partial:HH^2(\Gam,J)\xrightarrow{\sim}HH^3(\Gam,U)\).  Choose a normalised Hochschild \(2\)-cocycle \(\alpha:\Gam^{\otimes2}\to J\) representing \(\partial^{-1}\mu\) and define
\[
 A=\Gam\oplus J,\qquad (x,j)(y,l)=(xy,\ xl+jy+\alpha(x,y)).
\]
The cocycle identity is the associativity of this product, normalisation makes \((1,0)\) the unit, and \(J\) is a square-zero ideal.  Let \(M=I\) with \(A\) acting through \(A\to\Gam\), and \(\delta(i)=(0,q(i))\).  Then \(\delta\) is a bimodule map, both sides of \eqref{eq:peiffer} vanish because \(J\) annihilates \(I\), and \Cref{lem:crossed} gives a finite-dimensional two-term dg algebra \(D\) with \(H^0(D)=A/(0\oplus J)=\Gam\) and \(H^{-1}(D)=\ker q=U\).

For the class, use the section \(s(x)=(x,0)\) and a linear section \(j:J\to I\) of \(q\) with \(M'=j(J)\); then \(s(x)s(y)-s(xy)=(0,\alpha(x,y))=\delta(j\alpha(x,y))\), so \(\beta=j\alpha\) in \eqref{eq:beta}, and \Cref{lem:two-term-minimal} gives
\[
 m_3(x,y,z)=x\,j\alpha(y,z)-j\alpha(xy,z)+j\alpha(x,yz)-j\alpha(x,y)\,z .
\]
This is the Hochschild differential of the lift \(j\alpha\) of \(\alpha\) to \(I\); it takes values in \(U\) and represents \(\partial[\alpha]=\mu\) by the definition of the connecting homomorphism.

For the consequence, let \(B\) be proper with cohomology in \(\{-1,0\}\); if \(B\) is acyclic it is quasi-isomorphic to \(0\), realised by the identity of \(k\) and its contractible cone.  Otherwise the transfer theorem \cite[Section~3.3]{Keller01Ainf} provides a strictly unital minimal \(A_\infty\)-model on \(H^*(B)=\Gam\oplus U[1]\), whose structure is \((m_2,m_3)\) by \Cref{prop:higher-mult}; let \(\mu\) be the class of \(m_3\).  The dg algebra \(D\) just constructed has a minimal model with the same \(\Gam\), \(U\) and class, hence \(A_\infty\)-isomorphic to that of \(B\) by \Cref{lem:gauge}.  Two dg algebras with \(A_\infty\)-isomorphic minimal models are connected by a zigzag of quasi-isomorphisms of dg algebras \cite[Section~3.3]{Keller01Ainf}, \cite[Section~1.3]{LefevreHasegawa03}.  Apply \Cref{thm:realisation} to \(D\).
\end{proof}

\begin{corollary}[All profiles for a thin kernel]\label{cor:thin-all-profiles}
For every proper dg algebra \(B\) with cohomology concentrated in degrees \(-1\) and \(0\), the marked bounded defects realised by stably flat universal localisations of finite-dimensional algebras with kernel dg algebra \(B\) are exactly
\[
 \{\cU_E/\per B:E\in\Dfd(B)\}.
\]
All of them are idempotent complete and coincide with the corresponding singular defects before completion.
\end{corollary}

\begin{proof}
Necessity is \Cref{cor:pair-invariant}.  \Cref{thm:finite-crossed} provides a full-profile realisation whose target is projective as a left module over the source, and \Cref{thm:ambient}\ref{amb:d} realises every \(\cU_E\).  The free module of a connective model is a silting compact generator with finite-dimensional endomorphism algebra, so \Cref{cor:finite-dimensional-precompletion,cor:negative-k-completeness} apply.
\end{proof}

\begin{example}[A non-formal kernel]\label{ex:nonformal}
Let \(A\) have basis \(1,t,a,b,c\), with unit \(1\) and the only nonzero products of radical basis elements
\[
 t^2=c,\qquad ta=b,\qquad at=c,
\]
so that \(\rad^3A=0\).  Let \(M\) have basis \(z,l,r,w\), with the only nonzero radical actions
\[
 tz=l,\qquad tr=w,\qquad zt=r,\qquad lt=w,
\]
and let \(\delta z=a\), \(\delta l=b\), \(\delta r=c\), \(\delta w=0\).  The bimodule laws hold, the only nontrivial mixed check being \((tz)t=w=t(zt)\); \(\delta\) is bimodule-linear, and \eqref{eq:peiffer} holds because the image of \(\delta\) annihilates \(M\).  Hence \(D=[M\to A]\) is a nine-dimensional two-term dg algebra with
\[
 \Gam=k[t]/(t^2),\qquad U=kw,\qquad tU=Ut=0 .
\]
Since \(t^2=c=\delta(r)\), \Cref{lem:two-term-minimal} with \(s(t)=t\) and \(\beta(t,t)=r\) gives
\[
 m_3(t,t,t)=t\,r-r\,t=w\ne0 .
\]
For any normalised \(2\)-cochain \(f:\Gam^{\otimes2}\to U\) the Hochschild coboundary satisfies \((d_{\mathrm{Hoch}}f)(t,t,t)=tf(t,t)-f(t^2,t)+f(t,t^2)-f(t,t)t=0\), so \(m_3\) is not a coboundary: \(D\) is not formal, and \(HH^3(\Gam,U)\cong k\) is generated by its class.  \Cref{thm:realisation} realises \(D\) as the kernel of a stably flat universal localisation \(R_\delta\to M_2(A)\) at a single map, with \(\dim_kR_\delta=19\), \(\dim_kM_2(A)=20\), target projective on both sides and full bounded profile.  Its kernel dg algebra therefore carries a nonzero \(m_3\): kernel dg algebras of finite universal localisations of finite-dimensional algebras need not be formal.  The dg algebra, the ring, the two projective decompositions, the strict dg identification and the value of \(m_3\), computed once on \(D\) and once on the presentation-cone algebra, are checked exactly by \texttt{verify\_nonformal.py} (\Cref{app:computation}).
\end{example}

\begin{remark}[Finite strictification and the crossed model]\label{rem:finite-strict-priority}
The existence of a finite two-term model also follows from \cite[Appendix~A, Theorem~A.3]{Goodbody24}.  Take a finite non-positive strict model \(C\), and quotient by the dg ideal equal to \(C^i\) in degrees \(i\le-2\), to \(d(C^{-2})\) in degree \(-1\), and to zero in degree zero.  Nonpositivity and the Leibniz rule show that this is a dg ideal.  If \(H^{<-1}(C)=0\), the projection is a quasi-isomorphism onto a finite two-term dg algebra.  Thus finite-model existence is prior work; the construction in \Cref{thm:finite-crossed} explicitly starts from \((\Gam,U,\mu)\) and computes the resulting minimal-model class, before applying the ring-epimorphism realisation of \Cref{thm:realisation}.
\end{remark}

\subsection{Optimality of the t-exact realisation}

The construction above has a converse when the marked kernel embedding is required to preserve the standard t-structures.  The converse applies to arbitrary ordinary homological ring epimorphisms, without a projectivity assumption on the target.

\begin{theorem}[T-exact kernel realisations are two-term]\label{thm:texact-optimality}
Let \(\lambda:R\to S\) be a homological ring epimorphism of ordinary rings and suppose
\[
 \cK_\lambda=\ker(S\Ltensor_R-)=\Loc(T),
 \qquad T\in K^b(R\text{-}\mathrm{proj}).
\]
If \(B=\RHom_R(T,T)^{\op}\) is connective and its Morita embedding
\(F:D(B)\to D(R)\) is t-exact for the standard t-structures, then \(T\) is homotopy equivalent to a complex of finitely generated projectives concentrated in degrees \([-1,0]\).  In particular \(H^i(B)=0\) for \(i<-1\).
\end{theorem}

\begin{proof}
A fully faithful t-exact functor reflects standard cohomological bounds.  Indeed, applying it to a truncation triangle gives a truncation triangle, whence
\(H^n(FY)\cong F_{\heartsuit}(H^nY)\); the fully faithful heart functor reflects zero objects.  This argument does not require \(Y\) to be bounded.

Homologicality gives the localisation triangle
\[
 K_R\longrightarrow R\xrightarrow{\lambda}S\longrightarrow K_R[1],
 \qquad K_R=\Cone(R\to S)[-1]\in\cK_\lambda.
\]
To check the last assertion in all degrees, apply \(S\Ltensor_R-\): the resulting map \(S\to S\Ltensor_RS\) is the inverse of the multiplication isomorphism.  With our cohomological convention, the displayed representative of \(K_R\) has terms in degrees \(0,1\).  Thus, writing \(K_R\cong FC\), reflection gives \(C\in D^{[0,1]}(B)\).  For every \(n\), adjunction gives
\[
 \Hom_{D(R)}(T,S[n])
 \cong\Hom_{D(S)}(S\Ltensor_RT,S[n])=0.
\]
Applying \(\RHom_R(T,-)\) to the triangle identifies the cohomology of \(T^\vee=\RHom_R(T,R)\) with that of \(\RHom_R(T,K_R)\).  Full faithfulness and \(F(B)=T\) identify the latter groups with \(\Hom_{D(B)}(B,C[n])=H^n(C)\).  Hence
\[
 T^\vee\in D^{[0,1]}(R^{\op}),
 \qquad T=F(B)\in D^{\le0}(R).
\]
Only cohomological bounds are asserted here; the two complexes have different additional module actions.

For a bounded projective complex whose cohomology vanishes above degree \(b\), the differential onto its highest term above \(b\) is a split surjection.  Removing the associated contractible summand and repeating eliminates all terms above \(b\).  Apply this first to \(T\), obtaining a projective representative with no positive terms.  Its finite-projective dual has no negative terms.  Apply the same cancellation to that dual above degree one; it is then supported in \([0,1]\).  Finite-projective biduality gives a representative of \(T\) in \([-1,0]\).  All steps preserve homotopy equivalence.  Its endomorphism complex has terms only in degrees \([-1,1]\), proving the vanishing below \(-1\).
\end{proof}

\begin{corollary}[Universality of a t-exact realisation]\label{cor:texact-universal}
Under \Cref{thm:texact-optimality}, write a two-term representative as \(T=[P^{-1}\xrightarrow{\sigma}P^0]\).  Then \(\lambda\) is the universal localisation at \(\sigma\).
\end{corollary}

\begin{proof}
The vanishing \(S\Ltensor_RT=0\) says that \(S\otimes_R\sigma\) is invertible.  The essential image of derived restriction is \(\cK_\lambda^\perp\).  A stalk \(R\)-module \(M\) belongs to this image exactly when \(\Hom_R(\sigma,M)\) is invertible.  Its preimage is a stalk \(S\)-module, because restriction is exact and reflects cohomological vanishing.  Its extending \(S\)-action is unique: full faithfulness applied to the identity of the underlying \(R\)-module identifies any two such actions.

Let \(f:R\to C\) invert \(\sigma\).  Tensor--Hom adjunction puts \({}_RC\) in \(\cK_\lambda^\perp\), giving its unique extending \(S\)-action.  Every right multiplication on \(C\) is \(R\)-linear and therefore \(S\)-linear.  Consequently \(s\cdot c=(s\cdot1)c\), and \(g(s)=s\cdot1\) is a unital ring homomorphism \(S\to C\) extending \(f\).  Any other extension defines the same \(S\)-action, hence equals \(g\).  This is the required universal property.
\end{proof}

\begin{corollary}[Classification of t-exact marked realisations]\label{cor:texact-classification}
For a proper connective dg algebra \(B\) over a field, the following are equivalent:
\begin{enumerate}[label=\textup{(\roman*)}]
\item \(B\) is quasi-isomorphic to the dg endomorphism algebra, with opposite, of a perfect generator of the kernel of a homological epimorphism of finite-dimensional algebras, with standard t-exact marked kernel embedding;
\item \(H^i(B)=0\) for \(i<-1\);
\item such a realisation exists as a universal localisation at one map, with target projective on both sides and full bounded profile.
\end{enumerate}
\end{corollary}

\begin{proof}
The implication (i)\(\Rightarrow\)(ii) is \Cref{thm:texact-optimality}; (ii)\(\Rightarrow\)(iii) is \Cref{thm:finite-crossed,thm:realisation}; and (iii)\(\Rightarrow\)(i) is immediate.  Quasi-isomorphisms of connective dg models preserve the standard t-structures.
\end{proof}

Thus allowing ordinary homological epimorphisms which are not assumed to be universal localisations does not enlarge this t-exact realisation range.  The full-profile seed of \Cref{thm:finite-crossed} still realises every cut by \Cref{thm:ambient}; an extension realising a proper cut need not retain the seed's projectivity or t-exactness.

\subsection{Projective targets and the Schur boundary}

\begin{corollary}[Kernel categories with projective right target]\label{cor:projective-target-classes}
The dg Morita classes of compact kernel categories of homological epimorphisms of finite-dimensional algebras with projective right target are exactly the classes admitting a proper connective dg model with cohomology in \([-1,0]\).  Every such class has a realisation with target projective on both sides, standard t-exact marking by a thin model, and full bounded profile.
\end{corollary}

\begin{proof}
Let \(S_R\) be projective.  By \Cref{prop:left-projection}, the kernel has the perfect generator
\[
 G=\Cone\!\left(\Hom_{R^{\op}}(S,R)\longrightarrow R\right),
\]
with terms in \([-1,0]\) and \(\Hom(G,G[n])\cong H^n(G)\).  Thus \(\RHom_R(G,G)^{\op}\) is proper with cohomology in \([-1,0]\), and models the compact kernel.  Conversely apply \Cref{thm:finite-crossed,thm:realisation} to a thin model.  A dg Morita equivalence preserves \(\Dfd\), by its intrinsic finite-total-Hom characterisation, so the full-profile assertion is independent of the compact-generator marking.  The t-exact assertion concerns the thin marking supplied by the construction.
\end{proof}

\begin{proposition}[Negative Calabi--Yau obstruction]\label{prop:negative-cy-obstruction}
Let \(B\) be a proper dg algebra such that \(\per B\ne0\) has Serre functor \([-d]\), where \(d\ge2\).  Then \(\per B\) is not equivalent to the compact kernel of either
\begin{enumerate}[label=\textup{(\roman*)}]
\item a homological epimorphism of finite-dimensional algebras with \(S_R\in\per(R^{\op})\); or
\item a finite stably flat universal localisation of finite-dimensional algebras, without any perfection assumption on the target.
\end{enumerate}
Neither exclusion requires t-exactness or full profile.
\end{proposition}

\begin{proof}
Every nonzero \(G\in\per B\) satisfies
\[
 \Hom(G,G[-d])\cong D_k\End(G)\ne0
\]
by Serre duality.  In (i), \Cref{prop:left-projection} supplies a nonzero perfect generator with \(\Hom(G,G[n])=0\) for \(n<-1\), a contradiction.  In (ii), the transported sum of the presentation cones has the same vanishing by \Cref{prop:thin}, again a contradiction.
\end{proof}

\begin{example}[Two excluded formal algebras]\label{ex:negative-cy-algebras}
The dg algebras
\[
 B_2=k[u]/(u^2),\quad |u|=-2,
 \qquad
 B_{\wedge}=\Lambda_k(u,v),\quad |u|=|v|=-1,
\]
with zero differential satisfy \Cref{prop:negative-cy-obstruction} with \(d=2\), over every field.  For \(B_2\), take the trace with value one on \(u\) and zero on \(1\); for \(B_{\wedge}\), take the coefficient of \(uv\).  The Gram matrices of \((a,b)\mapsto\operatorname{tr}(ab)\), in the respective bases \((1,u)\) and \((1,u,v,uv)\), are
\[
 \begin{pmatrix}0&1\\1&0\end{pmatrix},
 \qquad
 \begin{pmatrix}0&0&0&1\\0&0&1&0\\0&-1&0&0\\1&0&0&0\end{pmatrix}.
\]
These pairings are nondegenerate and graded symmetric, including in characteristic two, and give dg bimodule isomorphisms \(D_kB\cong B[-2]\).  Finite-projective tensor--Hom duality gives
\[
 D_k\RHom_B(X,Y)\simeq
 \RHom_B\bigl(Y,D_kB\Ltensor_BX\bigr)
 \qquad(X,Y\in\per B).
\]
The formula holds on finite free modules and extends by shifts, cones and direct summands; it identifies \([-2]\) as the Serre functor.
\end{example}

\begin{example}[A Schur realisation outside the epimorphism range]\label{ex:sphere-schur}
Let \(A\) be the radical-square-zero algebra of the oriented three-cycle
\[
 1\xrightarrow{a_1}2\xrightarrow{a_2}3\xrightarrow{a_3}1.
\]
Its basis is \(e_1,e_2,e_3,a_1,a_2,a_3\), with orthogonal vertex idempotents, the source and target identities of the arrows, and all products of two arrows zero.  Put \(e=e_2+e_3\), so \(eAe=k(2\to3)\), and take
\[
 T=[Ae_2\xrightarrow{a_2}Ae_3\xrightarrow{a_3}Ae_1]
 \quad\text{in degrees }[-2,-1,0].
\]
Then there is a recollement
\[
 D(B_2)\longrightarrow D(A)\xrightarrow{e(-)}D(eAe)
\]
whose kernel embedding is t-exact and has full bounded profile.  The source has dimension six and the quotient corner has dimension three.  The quotient functor here is the Schur functor, not extension along an ordinary ring epimorphism.

Indeed, the endomorphism complex of \(T\) has dimensions \(1,0,3,2\) in degrees \(-2,-1,0,1\), and zero in every other degree.  Its only nonzero differential sends the three diagonal scalars to their two consecutive differences and has rank two over every field.  Thus its cohomology is \(k\) in degrees \(-2,0\).  Sending \(u\) to right multiplication by \(a_1\), from \(Ae_1\) to \(Ae_2\), and \(1\) to the identity gives a strict quasi-isomorphism
\[
 B_2\longrightarrow\End_A^\bullet(T)^{\op}.
\]
As right dg \(B_2\)-modules, \(e_1T=B_2\); the other two vertex components are identity-contractible two-term complexes, on which \(u\) acts trivially.  Consequently \(T\Ltensor_{B_2}-\) preserves cohomological degrees and finite-dimensional total cohomology, and reflects the latter.

The object \(T\) belongs to \(\ker e(-)\).  If \(X\) belongs to this kernel, the finite filtration of \(\Hom_A^\bullet(T,X)\) by the three terms of \(T\) has two acyclic associated factors, \(e_2X\) and \(e_3X\), leaving \(e_1X\).  Hence \(\RHom_A(T,X)=0\) implies \(X=0\).  The \(\Loc(T)\)-approximation triangle inside \(D(A)\) now gives \(\ker e(-)=\Loc(T)\): its residual object belongs to the kernel and is right-orthogonal to \(T\), so is zero.  Dg Morita theory and the idempotent recollement of \cite[Proposition~6.1]{KalckYangII} yield the displayed recollement.  The complete finite computation is included in the supplement.  By \Cref{thm:texact-optimality,prop:negative-cy-obstruction}, this example cannot be replaced by any of the ordinary-epimorphism realisations excluded there.
\end{example}

\begin{example}[Full profile does not force a thin marking]\label{ex:marking-boundary}
Let \(A=k(1\to2)\), \(R=A\times k\), and \(\lambda:R\to k\) be projection.  The target is projective on both sides and multiplication \(k\otimes_Rk\to k\) is an isomorphism, so all positive Tor groups vanish and the kernel is \(D(A)\).  Put \(P_i=Ae_i\) and \(T=P_1\oplus P_2[2]\).  The complete Hom calculation is
\[
 \End_A(P_i)=k,\qquad
 \Hom_A(P_1,P_2)=k,\qquad\Hom_A(P_2,P_1)=0;
\]
all higher or lower Hom groups of these stalk projectives vanish.  Therefore \(B=\End_R^\bullet(T)^{\op}\) has zero differential,
\[
 H^0(B)=k\times k,\qquad H^{-2}(B)=k,
\]
and no other cohomology.  Since \(T\) and \(A\) classically generate the same perfect category, the finite-total-Hom characterisation gives full profile for this marking.  It is not t-exact by \Cref{thm:texact-optimality}.  Thus \Cref{cor:projective-target-classes} concerns existence of a thin model in a dg Morita class; it does not assert that every connective marked model is thin.  In particular, the broader \Cref{q:realisation} is not settled by the t-exact classification.
\end{example}


\subsection{Open problems}

The standard t-exact version of \Cref{q:realisation} is settled by \Cref{cor:texact-classification}, and \Cref{prop:negative-cy-obstruction} excludes a class of dg Morita types even without t-exactness.  The general full-profile question remains open, as does realisation by unrestricted ordinary homological epimorphisms after dropping t-exactness and allowing infinite right target projective dimension.  The Schur construction in \Cref{ex:sphere-schur} is a different realisation problem.

\begin{question}\label{q:profile-m3}
For \Cref{ex:nonformal}, do the posets of far-away orthogonals \(\{\cU_E/\per B\}\) of the non-formal dg algebra and of the formal dg algebra with the same cohomology differ?  The two have different Ext-algebras of the residue field, but eventual orthogonality is not determined by low-degree Ext.
\end{question}

\begin{question}\label{q:nonsilting}
The realisation theorem concerns kernels whose cone generator is silting.  Which dg algebras with cohomology in degrees \(-1,0,1\) occur as the endomorphism dg algebra of a presentation cone of a stably flat universal localisation of finite-dimensional algebras, as in the example \(k\times k\) after \Cref{prop:thin}?
\end{question}
